\section{O Mecanismo de Penalização e Atenuação}
\label{sec:penalization}


\subsection{Cenários e a Função de Penalização}
\label{subsec:scenarios}

Nesta etapa, a rede de deslocamentos é avaliada em dois estados distintos: (1) Cenário Típico, que reflete as condições normais de mobilidade; e (2) Cenário Condicionado, no qual os custos da rede são alterados por variáveis ambientais.

No \textbf{Cenário Típico}, o peso de cada aresta do grafo ($W_{base}$) corresponde estritamente ao comprimento real do trecho da via. O tempo de percurso em condições normais é obtido dividindo-se essa distância pela velocidade média de deslocamento definida para o modo de transporte analisado.

No \textbf{Cenário Condicionado}, os custos de deslocamento são recalculados por meio de uma função de penalização. Esta função cruza a malha viária com as camadas de condicionantes ambientais ($C$), gerando um peso penalizado ($W_{condicionado}$) para cada aresta, estruturado na forma $W_{condicionado} = f(W_{base}, C)$. A aplicação geométrica dessa função varia conforme a natureza do dado ambiental:

\begin{itemize}
    \item \textbf{Dados contínuos} (ex.: \textit{rasters} de temperatura): O algoritmo extrai a média dos valores dos pixels que se sobrepõem geograficamente à aresta. Esse valor atua como um penalizador contínuo, aumentando o tempo de percurso proporcionalmente à severidade da condição ambiental ao longo daquele trecho.
    
    \item \textbf{Dados discretos} (ex.: polígonos de alagamento ou barreiras físicas): O algoritmo utiliza operações de interseção vetorial. A função pode aplicar multiplicadores severos ou atribuir peso infinito aos trechos interceptados, simulando interdições totais que forçam o roteamento a buscar caminhos alternativos. Essa mesma lógica de interseção pode ser usada de forma inversa, reduzindo pesos para simular áreas favoráveis (como vias arborizadas).
\end{itemize}

A forma matemática da função de penalização $f$ (linear, exponencial, limiar de corte) é um parâmetro ajustável. A escolha e calibração desses parâmetros devem considerar dados ou estudos do contexto local, já que a influência de variáveis ambientais sobre a velocidade e o esforço para caminhar varia conforme o clima, a cultura e as características urbanas. Evidências recentes mostram que fatores como temperatura, clima e condições viárias alteram significativamente o padrão e a velocidade da caminhada \cite{LIANG2020106811,Obuchi2021}. Por isso, recomenda-se que o aplicador utilize valores observados localmente ou reportados na literatura para definir as funções de penalização e os parâmetros do modelo.

Por fim, com os grafos dos dois cenários estruturados, um algoritmo de caminho de menor custo é executado para calcular as matrizes de tempo de viagem. A partir desses tempos, computam-se os indicadores apresentados na Seção \ref{subsubsec:indicadores}: o PTh de cada unidade territorial, o índice G (desigualdade) e o F15 (fração da população atendida) para ambos os cenários.

\subsection{Análise Comparativa}
\label{subsec:comparative}

Nesta etapa, quantificam-se as perdas de acessibilidade comparando diretamente os indicadores do cenário típico com os do cenário condicionado. O objetivo é mensurar a degradação da mobilidade imposta pelos fatores ambientais e identificar pontos críticos onde a dificuldade de circulação provoca maior isolamento.

Essa comparação é estruturada em três procedimentos analíticos principais:

\begin{itemize}
    \item \textbf{Variação relativa:} Cálculo das diferenças absolutas e percentuais dos tempos de deslocamento e dos níveis de acesso entre os dois cenários, evidenciando a magnitude da perda para cada unidade da malha.
    \item \textbf{Validação estatística:} Aplicação de testes de hipótese para confirmar a significância das alterações nos tempos de acesso provocadas pela perturbação ambiental em relação à rede original.
    \item \textbf{Mapeamento espacial:} Uso de métodos de associação local ou clusterização sobre as diferenças de tempo. Isso permite identificar, agrupar e mapear as manchas de degradação, revelando geograficamente as áreas mais penalizadas da cidade.
\end{itemize}

\subsection{Análise de Desigualdade Socioeconômica}
\label{subsec:socioeconomical}

Nesta etapa, as perdas de acessibilidade são analisadas em conjunto com os atributos socioeconômicos da população, com o objetivo de identificar padrões de vulnerabilidade associados às perturbações ambientais.

Cada célula da malha de análise, já associada aos dados sociodemográficos por compatibilização espacial, é tratada como unidade de observação. Para cada célula, consideram-se tanto os indicadores de acessibilidade no cenário típico ($PT_h$) quanto as variações induzidas pela perturbação ($\Delta PT_h$).

Inicialmente, aplica-se um modelo de regressão linear, no qual $PT_h$ e $\Delta PT_h$ são utilizados como variáveis dependentes em especificações distintas. As variáveis independentes incluem renda média, nível de escolaridade, densidade populacional e demais indicadores disponíveis. Esse procedimento permite estimar a associação entre características socioeconômicas e níveis de acessibilidade, bem como sua variação sob condições adversas.

Complementarmente, realiza-se uma análise de agrupamento socio-espacial, na qual as células são agrupadas considerando simultaneamente sua proximidade geográfica e a similaridade em termos de atributos socioeconômicos e indicadores de acessibilidade. Esse procedimento resulta em regiões adjacentes e internamente homogêneas, permitindo identificar padrões territoriais de desigualdade e vulnerabilidade de forma integrada.

% https://geoportal.sgb.gov.br/geosgb/

\subsection{Limitações}
\label{subsec:limitations}

A calibração dos efeitos ambientais no modelo é limitada pela disponibilidade e qualidade dos dados ou por falta de estudos específicos. Por isso, os valores de indicadores como o PTh podem não representar tempos absolutos reais, sendo mais adequados para análises comparativas entre diferentes regiões do que para estimativas precisas de tempo. Dessa forma, o modelo cumpre melhor o papel de identificar padrões de desigualdade e vulnerabilidade territorial diante de perturbações ambientais do que de quantificar o tempo exato de acesso.

Além disso, há limitações relacionadas ao comportamento dos indivíduos. Em situações extremas, as pessoas podem optar por não se deslocar ou escolher modos de transporte diferentes da mobilidade ativa. O modelo pressupõe a manutenção do deslocamento a pé mesmo sob condições adversas, o que pode não refletir a realidade observada. Assim, os resultados simulam apenas o impacto potencial das perturbações na acessibilidade da rede, não substituindo análises que considerem efetivamente as escolhas modais dos usuários. Ressalta-se, ainda, que a metodologia pode ser adaptada para outros modos de transporte, mediante ajustes nos parâmetros de velocidade e custo.

Ainda, os indicadores calculados refletem apenas o potencial de acesso espacial e o esforço exigido para atingir determinados destinos, sem considerar fatores como a qualidade, capacidade, funcionamento ou atratividade dos serviços presentes nos POIs. Além disso, esses indicadores não capturam obstáculos subjetivos, como percepção de segurança, conforto ou preferência individual de rotas. Por isso, apesar de úteis para análise territorial comparativa, os resultados não devem ser interpretados como medida direta da experiência do usuário ou garantia de acesso efetivo aos serviços analisados.
